% !TeX encoding = UTF-8
% maindoc



\documentclass[a4paper,12pt,oneside,german,toc=bibliography]{scrbook} 


\usepackage[ngerman]{babel}
\usepackage[utf8]{inputenc}
\usepackage{amsmath,amsthm,amssymb,amsfonts,amscd,amsbsy,amsxtra,amsthm}

\usepackage[alphabetic,nobysame,bibtex-style]{amsrefs}
\makeatletter
\renewcommand\PrintNames@a[4]{%
    \PrintSeries{\name}
    {#1}
    {}{ und \set@othername}
    {;}{ \set@othername}
    {}{ und \set@othername}
    {#2}{#4}{#3}%
}
\makeatother


\usepackage{dsfont}
\usepackage{enumerate}
\usepackage{graphicx}
\usepackage{subfig}
\usepackage{upgreek,textcomp}
\usepackage{cleveref}
\usepackage[Algorithmus]{algorithm}
\usepackage{algpseudocode}
\usepackage[autostyle=true,german=quotes]{csquotes}
\usepackage{ziffer}
\usepackage{booktabs}

\makeatletter
\newcommand{\specialcell}[1]{\ifmeasuring@#1\else\omit$\displaystyle#1$\ignorespaces\fi}
\makeatother

\usepackage[numbered,autolinebreaks,bw]{mcode}   %framed

\theoremstyle{definition}
\newtheorem{defn}{Definition}[section]
\newtheorem{bem}[defn]{Bemerkung}

\theoremstyle{plain}
\newtheorem{satz}[defn]{Satz}
\newtheorem{lem}[defn]{Lemma}
\newtheorem{prop}[defn]{Proposition}
\newtheorem{kor}[defn]{Korollar}

\numberwithin{equation}{section}

\DeclareMathOperator{\cond}{cond}
\DeclareMathOperator{\spn}{span}
\DeclareMathOperator{\diag}{diag}


\newcommand{\R}{\mathds{R}}
\newcommand{\C}{\mathds{C}}
\newcommand{\N}{\mathds{N}}
\newcommand{\Q}{\mathds{Q}}
\newcommand{\la}{\lambda}
\newcommand{\calV}{\mathcal{V}}
\newcommand{\calW}{\mathcal{W}}
\newcommand{\calP}{\mathcal{P}}
\newcommand{\calO}{\mathcal{O}}
\newcommand{\mPhi}{\mathit{\Phi}}
\newcommand{\eqItem}{\item~\vspace{-2\normalbaselineskip}}
\newcommand{\imaginary}{\mathrm{i}}
\newcommand{\LskalProd}{(\,\cdot\,,\,\cdot\,)_w}

\let\vphi=\phi
\renewcommand{\phi}{\varphi}

\newcommand{\transp}[1]{#1^\mathrm{T}}


\begin{document}

%!TEX root = /Users/mbantle/Dropbox/BantleM/Thesis/diplomarbeit.tex

\begin{titlepage}
	\begin{figure}[h]
		\includegraphics[width=\textwidth]{images/unilogoA4}
	\end{figure}
	\vspace{1cm}
\begin{center}
{\large \bfseries Fakultät für Mathematik und Wirtschaftswissenschaften \\
\vspace{0.3cm}
\bfseries Institut für Numerische Mathematik \\
}
\vspace{2cm}
{\Large
Projekt CSE\\}

\vspace{2cm}
{\LARGE\bfseries
Radarsignalverarbeitung mit Grafikprozessoren\\
}
\end{center}

%
\vspace{3.5cm}
%
\begin{center}
%\vspace{0.2cm}
%{\bfseries\mbox{Anton Hügel}} \\
%\vspace{0.2cm}
%Zwischenstand vom \mbox{\today} \\

vorgelegt von \\
\vspace{0.2cm}
{\bfseries\mbox{Anton Hügel, Jonas Schwer, Lukas Tatzel, Michael Thoma}} \\
\vspace{0.2cm}
am \mbox{\today} \\

\vspace*{1.2cm}
{\bfseries Betreuung} \\
\mbox{}\\
\mbox{Prof. Dr. Stefan A. Funken, Dr. Markus Bantle} \\
\end{center}
\vfill
\end{titlepage}


\tableofcontents

%%%%%%%%%%%%%%%%%%%%%%%%%%%%%%%%%%%%%%%%%%%%%%%%%
%%%%%%%%%%%%%%%%%%%%%%%%%%%%%%%%%%%%%%%%%%%%%%%%%

\chapter{Theorie}
In diesem Kapitel soll die Funktionsweise eines Radar-Systems (\textbf{ra}dio \textbf{d}etection \textbf{a}nd \textbf{r}anging) und dessen Signalverarbeitungsalgorithmen sowie einige mathematische Grundlagen kurz erläutert werden, wie sie ausführlicher auch in \cite{Richards,RSH,Ludloff} nachzulesen sind. Zudem wird eine grundlegende Einführung in die Programmierschnittstelle OpenCL gegeben, mit welcher eine Programmierung für Grafikprozessoren umgesetzt werden kann.
%%%%%%%%%%%%%%%%%%%%%%%%%%%%%%%%%%%%%%%%%%%%%%%%%

\section{Funktionsweise eines Radars} 
    % Aufbau
\subsection{Aufbau eines Radar}

    %Puls-Doppler-Radar

\subsection{Puls-Doppler-Radar}

Beim Puls-Doppler-Radar werden durch die Antenne des Radars kurze Sendepulse in dieselbe Richtung emittiert. Auf jeden Sendepuls folgt eine gewissen Empfangszeit, in der die Antenne auf Empfang geschaltet wird. Ein Vorteil des Puls-Doppler-Radars besteht also in der Möglichkeit, eine einzige Antenne sowohl für das Senden als auch das Empfangen verwenden zu können. 

Vergleicht man einen Sendepuls mit dessen Echo, so stellt man bei bewegten Objekten einen Unterschied im Frequenzverlauf fest. Dieses Phänomen wird als \glqq Dopplereffekt\grqq bezeichnet. Durch den Vergleich von mehreren Echos verschiedener Sendepulse lassen sich Rückschlüsse auf die Geschwindigkeit (genauer die Radialgeschwindigkeit bezogen auf das Radar) von detektierten Objekten ziehen. 
    %Signale

\subsection{Sende- und Empfangssignale}

%%%%%%%%%%%%%%%%%%%%%%%%%%%%%%%%%%%%%%%%%%%%%%%%%%

\section{Radarsignalverarbeitung}
    %Pulskompression

\subsection{Pulskompression}
    %Dopplerfilterung

\subsection{Dopplerfilterung}
    %Betrags-Bildung

\subsection{Betrags-Bildung}
    %CFAR

\subsection{CFAR}


%%%%%%%%%%%%%%%%%%%%%%%%%%%%%%%%%%%%%%%%%%%%%%%%%%%%%%%%%%%%%%%%%%%%%%%%%

\section{Fouriertransformation}
    %KFT

\subsection{Kontinuierliche FT}
    %DFT

\subsection{Diskrete FT (DFT)}
    %FFT

\subsection{Fast-Fourier-Transform (FFT)}



%%%%%%%%%%%%%%%%%%%%%%%%%%%%%%%%%%%%%%%%%%%%%%%%%%%%%%%%%%%%%%%%%%%%%%%%%

\section{OpenCL}


%%%%%%%%%%%%%%%%%%%%%%%%%%%%%%%%%%%%%%%%%%%%%%%%%%%%%%%%%%%%%%%%%%%%%%%%%
%%%%%%%%%%%%%%%%%%%%%%%%%%%%%%%%%%%%%%%%%%%%%%%%%%%%%%%%%%%%%%%%%%%%%%%%%


\chapter{Projekt}
Eine Abhandlung des bearbeiteten Projekts wird in diesem Kapitel gegeben,
wobei zunächst die Aufgagebenstellung dargelegt, dann deren Umsetzung erläutert
und zuletzt eine Verifikation und Bewertung der Ergebnise durchgeführt wird.

%%%%%%%%%%%%%%%%%%%%%%%%%%%%%%%%%%%%%%%%%%%%%%%%%%%%%%%%%%%%%%%%%%%%%%%%%

\section{Anforderungen}

%%%%%%%%%%%%%%%%%%%%%%%%%%%%%%%%%%%%%%%%%%%%%%%%%%%%%%%%%%%%%%%%%%%%%%%%%

\section{Implementierung}

%%%%%%%%%%%%%%%%%%%%%%%%%%%%%%%%%%%%%%%%%%%%%%%%%%%%%%%%%%%%%%%%%%%%%%%%%

\section{Verifikation}
    \subsection{Tests}     
    \subsection{Benchmarks}

%%%%%%%%%%%%%%%%%%%%%%%%%%%%%%%%%%%%%%%%%%%%%%%%%%%%%%%%%%%%%%%%%%%%%%%%%

\section{Zusammenfassung und Fazit}


%%%%%%%%%%%%%%%%%%%%%%%%%%%%%%%%%%%%%%%%%%%%%%%%%%%%%%%%%%%%%%%%%%%%%%%%%
%%%%%%%%%%%%%%%%%%%%%%%%%%%%%%%%%%%%%%%%%%%%%%%%%%%%%%%%%%%%%%%%%%%%%%%%%


\appendix

%%%%%%%%%%%%%%%%%%%%%%%%%%%%%%%%%%%%%%%%%%%%%%%%%%%%%%%%%%%%%%%%%%%%%%%%%

\chapter{Algorithmen}

%%%%%%%%%%%%%%%%%%%%%%%%%%%%%%%%%%%%%%%%%%%%%%%%%%%%%%%%%%%%%%%%%%%%%%%%%

\chapter{Quellcode}


%%%%%%%%%%%%%%%%%%%%%%%%%%%%%%%%%%%%%%%%%%%%%%%%%%%%%%%%%%%%%%%%%%%%%%%%%
%%%%%%%%%%%%%%%%%%%%%%%%%%%%%%%%%%%%%%%%%%%%%%%%%%%%%%%%%%%%%%%%%%%%%%%%%


% Literatur

\begin{bibdiv}
 %   \addcontentsline{toc}{chapter}{Literaturverzeichnis}
\begin{biblist}

\bib{RSH}{book} {
    author = {{Richards,}, Mark A.}*{inverted={yes}},
    author = {{Scheer,}, James A.}*{inverted={yes}},
    author = {{Holm,}, William A.}*{inverted={yes}},
    title = {Principles of modern radar},
    subtitle = {Basic Principles},
    address = {Raleigh},
    year = {2010},
    publisher = {Scitech Publishing},
    edition = {1. Auflage}
}

\bib{Richards}{book} {
    author = {{Richards,}, Mark A.}*{inverted={yes}},
    title = {Fundamentals of Radar Signal Processing},
    address = {Raleigh},
    year = {2005},
    publisher = {McGraw-Hill},
    edition = {1. Auflage}
}

\bib{Ludloff}{book} {
    author = {{Ludloff,}, Albrecht}*{inverted={yes}},
    title = {Praxiswissen Radar und Radarsignalverarbeitung},
    address = {Braunschweig/Wiesbaden},
    year = {2002},
    publisher = {Vieweg},
    edition = {3. überarbeitete Auflage}
}


\bib{Atkinson}{book} {
    author = {{Atkinson,}, Kendall}*{inverted={yes}},
    author = {{Han,}, Weinmin}*{inverted={yes}},
    title = {Theoretical Numerical Analysis},
    subtitle = {A Functional Analysis Framework},
    address = {NewYork},
    year = {2001},
    publisher = {Springer},
    edition = {1. Auflage}
}

\bib{MatheBibel}{book} {
    author = {{Arens,}, Tilo}*{inverted={yes}},
    author = {{Hettlich,}, Frank}*{inverted={yes}},
    author = {{Karpfinger,}, Christian}*{inverted={yes}},
    author = {{Kockelkorn,}, Ulrich}*{inverted={yes}},
    author = {{Lichtenegger,}, Klaus}*{inverted={yes}},
    author = {{Stachel,}, Hellmuth}*{inverted={yes}},
    title = {Mathematik},
    address = {Heidelberg},
    year = {2010},
    publisher = {Spektrum Akademischer Verlag},
    edition = {2., korrigierter Nachdruck}
}

\bib {Urban}{book} {
    author = {{Arendt,}, Wolfgang}*{inverted={yes}},
    author = {{Urban,}, Karsten}*{inverted={yes}},
    title = {Partielle Differenzialgleichungen: eine Einführung in analytische und numerische Methoden},
    address = {Heidelberg},
    year = {2010},
    publisher = {Spektrum Akademischer Verlag}
}

\bib{MBantle_Diss}{book}{
    title={On hp-Boundary ElementMethods for the Laplace Operator in Two Dimensions},
    author={{Bantle,}, Markus}*{inverted={yes}},
    organization = {Dissertation, Universität Ulm},
    date={2013}
}

\bib {Bunse} {book} {
    author = {{Bunse,}, Wolfgang}*{inverted={yes}},
    author = {{Bunse-Gerstner,}, Angelika}*{inverted={yes}},
    title = {Numerische lineare Algebra},
    address = {Stuttgart},
    year = {1985},
    publisher = {Teubner}
}

\bib{epsBEM}{book}{
    year = {2013},
    author = {{Bantle,}, Andreas}*{inverted={yes}},
    author = {{Bantle,}, Markus}*{inverted={yes}},
    author = {{Funken,}, Stefan}*{inverted={yes}},
    title = {epsBEM, efficient p-stable Matlab implementation of 2d BEM for Laplace and Lamé problems},
    publisher = {Technischer Bericht, Universität Ulm}
}

\bib {Hanke}{book} {
    author = {{Hanke-Bourgeois,}, Martin}*{inverted={yes}},
    title = {Grundlagen der Numerischen Mathematik und des Wissenschaftlichen Rechnens},
    address = {Wiesbaden},
    year = {2009},
    publisher = {Vieweg + Teubner},
    edition = {3., aktualisierte Auflage}
}

\bib{matlab}{book}{
    year = {2014},
    author = {MATLAB},
    title = {Version 8.4.0.150421 (R2014b)},
    publisher = {Software und Dokumentation, The MathWorks Inc.},
    address = {Natick, Massachusetts}
}

\bib{AMeister}{book}{
    author = {{Meister,}, Andreas}*{inverted={yes}},
    title = {Numerik linearer Gleichungssysteme: Eine Einführung in moderne Verfahren},
    address = {Wiesbaden},
    year = {2015},
    publisher = {Springer Spektrum},
    edition = {5., überarbeitete Auflage}
}

\bib{QSS}{book} {
    author = {{Quarteroni,}, Alfio}*{inverted={yes}},
    author = {{Sacco,}, Riccardo}*{inverted={yes}},
    author = {{Saleri,}, Fausto}*{inverted={yes}},
    title = {Numerische Mathematik. - 1.},
    address = {Berlin, Heidelberg},
    year = {2002},
    publisher = {Springer},
    pages = {XIV, 367 S.}
}

\bib{QSS2}{book} {
    author = {{Quarteroni,}, Alfio}*{inverted={yes}},
    author = {{Sacco,}, Riccardo}*{inverted={yes}},
    author = {{Saleri,}, Fausto}*{inverted={yes}},
    title = {Numerische Mathematik. - 2.},
    address = {Berlin, Heidelberg},
    year = {2002},
    publisher = {Springer}
}

\bib{Schwarz}{book}{
    author = {{Schwarz,}, Hans Rudolf}*{inverted={yes}},
    author = {{Köckler,}, Norbert}*{inverted={yes}},
    title = {Numerische Mathematik},
    address = {Wiesbaden},
    year = {2006},
    publisher = {Teubner},
    edition = {6., überarbeitete Auflage}
}

\bib{Olaf}{book} {
    author = {{Steinbach,}, Olaf}*{inverted={yes}},
    title = {Numerische Näherungsverfahren für elliptische Randwertprobleme: finite Elemente und Randelemente},
    address = {Wiesbaden},
    year = {2003},
    publisher = {Teubner},
    edition = {1. Auflage}
}

\bib{JWerner}{book} {
    author = {{Werner,}, Jochen}*{inverted={yes}},
    title = {Lineare und nichtlineare Gleichungssysteme, Interpolation, numerische Integration},
    address = {Braunschweig, Wiesbaden},
    year = {1992},
    publisher = {Vieweg}
}

\end{biblist}
\end{bibdiv}



\end{document}
